\section{Excursus: Octonions and exceptional geometry (non-essential)}
\label{sec:octonion}
%=============================================================================

	extbf{Purpose.} Earlier drafts experimented with exceptional-algebra language
(octonions) as a heuristic lens for certain geometric estimates.
This material is \emph{not used} in any rigorous step of the proof program and is
kept only as historical context.

	extbf{Warning.} A shared appearance of the number ``8'' (e.g. $\dim SU(3)=8$ and
$\dim_{\mathbb{R}}\mathbb{O}=8$) does not by itself imply any analytic advantage for
four-dimensional Yang--Mills measure estimates. Any functional-inequality or
spectral bound must be proved directly for the lattice measure.

\subsection{A correct geometric remark}

It is true that $G_2$ is the automorphism group of the octonions and that there
are classical relationships between $G_2$-geometry and certain homogeneous spaces.
However, these relationships do not yield, by themselves, rigorous improvements
for non-Abelian lattice gauge functional inequalities.

\subsection{Removed claims}

Previous drafts stated (i) an ``octonion-enhanced'' curvature mechanism and
(ii) an ``enhanced'' log-Sobolev inequality for the full Yang--Mills lattice
measure. These statements were not supported by a complete proof and are
therefore removed from the rigorous narrative.

\subsection{Status}

No octonion-based improvement is claimed.
Any group-specific constants used in functional inequalities must be derived by
standard analytic techniques (e.g. Bakry--\'Emery on compact groups, Holley--Stroock
perturbations where applicable) and the resulting bounds must be stated with all
assumptions explicit.

%=============================================================================






